\section{Phase 2a: Zerocheck Integration}
\label{sec:zerocheck-phase2a}

Phase 2a delivers the sumcheck infrastructure that replaces the quotient
polynomial in the WHIR fast path.  The integration is complete and
end-to-end verified for fibonacci; the remaining soundness gap (transition
constraints not yet bound to a PCS opening) is scheduled for the follow-up
phase.

\subsection{What Ships}

\paragraph{Sumcheck core (\texttt{crates/stark/src/zerocheck\_prover.rs}, 5 unit tests, all passing).}
\begin{itemize}
  \item \texttt{eq\_mle\_table(r)} -- dense evaluation of $\mathsf{eq}(r,\cdot)$ over $\{0,1\}^m$.
  \item \texttt{MultilinearExt::\{fold\_first, evaluate\}} -- standard MLE operations.
  \item \texttt{prove\_zerocheck\_multilinear(c\_evals, r)} -- deterministic prover.
  \item \texttt{verify\_zerocheck\_rounds(proof, r, initial\_claim)} -- deterministic verifier.
  \item \texttt{prove\_zerocheck\_with\_challenger(c\_evals, m, challenger)} -- Fiat-Shamir prover.
  \item \texttt{verify\_zerocheck\_with\_challenger(proof, m, initial\_claim, challenger)} -- Fiat-Shamir verifier.
\end{itemize}

\paragraph{Hypercube constraint evaluator.}
\texttt{eval\_constraints\_on\_hypercube} iterates each row of the Boolean
hypercube, lifts the (base-field) main and preprocessed rows to the
extension field, and invokes the chip's \texttt{Air::eval} via
\texttt{VerifierConstraintFolder}.  Produces a length-$2^m$ table of
extension-field values representing
$$C_{\mathrm{batched}}(b) = \sum_j \alpha^j \cdot C_j\bigl(\mathrm{row}_b, \mathrm{row}_{b+1}, \mathrm{preproc}_b, \mathrm{preproc}_{b+1}, \mathrm{pub}\bigr).$$

Selectors are set analytically per row ($\texttt{is\_first\_row}$,
$\texttt{is\_last\_row}$, $\texttt{is\_transition}$); next-row accesses use
cyclic wrap within the hypercube.

\paragraph{\texttt{ShardProof::zerocheck\_proofs} field.}
New optional field serialized with \texttt{\#[serde(default)]} for
backward-compatibility:
\begin{verbatim}
pub struct ShardProof<SC: StarkGenericConfig> {
    // existing fields
    #[serde(default)]
    pub zerocheck_proofs: Option<Vec<ZerocheckProof<Challenge<SC>>>>,
}
\end{verbatim}

\paragraph{Prover wiring (WHIR fast path).}
After sampling $\alpha$:
\begin{itemize}
  \item If \texttt{chip.permutation\_width() > 0} (lookup-dependent chips):
        emit an empty placeholder proof.  These chips become eligible for
        zerocheck once Phase~2b removes the permutation argument.
  \item Otherwise: run \texttt{eval\_constraints\_on\_hypercube} and
        \texttt{prove\_zerocheck\_with\_challenger}.
\end{itemize}
Timing is logged to \texttt{/tmp/ziren\_open\_breakdown.txt} under
\texttt{ZEROCHECK total=...ms}.

\paragraph{Verifier wiring.}
After sampling $\alpha$, if \texttt{zerocheck\_proofs} is present, for each
chip/proof pair the verifier replays the sumcheck transcript.  Empty
\texttt{rounds} means "skipped" (permutation-dependent chip).  Non-empty
proofs are passed to \texttt{verify\_zerocheck\_with\_challenger}.  New
error variants: \texttt{ZerocheckFailed}, \texttt{InvalidProofShape}.

\paragraph{Downstream fixes.}
\texttt{ShardProof} literals updated in
\texttt{crates/sdk/src/provers/mock.rs} and
\texttt{crates/recursion/circuit/src/stark.rs}.

\subsection{End-to-End Test Results}

\begin{table}[h]
\centering
\begin{tabular}{lll}
\toprule
\textbf{Mode} & \textbf{Result} & \textbf{Notes} \\
\midrule
FRI (quotient + permutation)                            & PASS (8.1\,s) & Existing pipeline unchanged \\
WHIR core (zerocheck emitted, skipped for all 19 chips) & PASS (6.8\,s) & Serialization round-trip works \\
WHIR compressed                                          & PASS (6.9\,s) & Recursion verifier bypasses empty proofs \\
\bottomrule
\end{tabular}
\caption{Phase 2a end-to-end verification on fibonacci.}
\label{tab:phase2a-e2e}
\end{table}

On fibonacci, all 19 CPU-program chips have \texttt{permutation\_width() > 0},
so \texttt{ZEROCHECK total=0ms chips=19} --- the sumcheck is never actually
executed for this workload.  Real-world exercise of the zerocheck math
will come once Phase~2b removes the permutation argument.

\subsection{Phase 2a Scope Limitation (By Design)}

Chips that use the lookup permutation argument are skipped because their
\texttt{Air::eval} calls \texttt{builder.permutation()}, which requires a
permutation trace we don't have in the WHIR fast path.  Two ways to lift
this:

\begin{enumerate}
  \item \textbf{Phase 2b} (LogUp-GKR): replace the permutation argument
        with GKR.  Chips no longer need \texttt{builder.permutation()}.
        All chips become zerocheck-eligible.
  \item \textbf{Interim fallback}: generate the permutation trace in the
        WHIR fast path purely for the zerocheck eval (not committed).
\end{enumerate}

Phase~2b is the correct long-term fix and matches SP1's hypercube
architecture.

\subsection{Files Changed (Phase 2a)}

\begin{verbatim}
crates/stark/src/lib.rs                    +1  module declaration
crates/stark/src/zerocheck.rs              +1  derive Serialize/Deserialize
crates/stark/src/zerocheck_prover.rs       +420 NEW
crates/stark/src/types.rs                  +5   ShardProof field
crates/stark/src/prover.rs                 +55  WHIR path wiring
crates/stark/src/verifier.rs               +40  zerocheck verify + errors
crates/sdk/src/provers/mock.rs             +1
crates/recursion/circuit/src/stark.rs      +1
\end{verbatim}

\section{Phase 2b: LogUp-GKR Protocol Core}
\label{sec:logup-gkr-phase2b}

Phase 2b Step 1 delivers a complete, unit-tested LogUp-GKR fraction-sum
protocol in \texttt{crates/stark/src/logup\_gkr.rs}.  The protocol takes a
set of leaf fractions $\{(\text{num}_i, \text{denom}_i)\}$ and lets the
prover commit to the claim
$$\sum_i \frac{\text{num}_i}{\text{denom}_i} = 0$$
via a logarithmic-depth reduction.  Each GKR layer is a \emph{folding}
step (not a sumcheck round) because the fraction-combine identity is
bilinear in the two halves.

\subsection{Fraction-Sum Combine}

For two fractions $(a, b) = a/b$ and $(c, d) = c/d$,
$$
\frac{a}{b} + \frac{c}{d} = \frac{a \cdot d + b \cdot c}{b \cdot d}.
$$
We encode this as the binary operation
$$(a,b) \oplus (c,d) = (a \cdot d + b \cdot c,\; b \cdot d).$$

The prover builds a balanced binary tree over leaf fractions using
$\oplus$; the root is $(N_m, D_m)$.  For a balanced LogUp argument,
$N_m = 0$, which the verifier checks in the clear.

\subsection{GKR Reduction}

At each layer $k$ (from root $m$ down to leaves $0$), the current claim
is a multilinear evaluation of $(N_k, D_k)$ at a point $z \in \mathbb{F}_{\mathrm{ext}}^{m-k}$.
The identity
\begin{align*}
N_{k-1}(z \Vert 0) \cdot D_{k-1}(z \Vert 1) + D_{k-1}(z \Vert 0) \cdot N_{k-1}(z \Vert 1) &= N_k(z), \\
D_{k-1}(z \Vert 0) \cdot D_{k-1}(z \Vert 1) &= D_k(z),
\end{align*}
lets the prover produce four values $(N_k(z\Vert 0), N_k(z\Vert 1), D_k(z\Vert 0), D_k(z\Vert 1))$
that the verifier checks against the current claim via the equations
above.  The verifier samples $r^* \in \mathbb{F}_{\mathrm{ext}}$ and
folds the claim to the new point $z \Vert r^*$:
\begin{align*}
N_{k-1}(z \Vert r^*) &= (1 - r^*) \cdot N_{k-1}(z \Vert 0) + r^* \cdot N_{k-1}(z \Vert 1), \\
D_{k-1}(z \Vert r^*) &= (1 - r^*) \cdot D_{k-1}(z \Vert 0) + r^* \cdot D_{k-1}(z \Vert 1).
\end{align*}

After $m$ layers, the remaining claim is
$\bigl(N_0(\mathsf{eval\_point}), D_0(\mathsf{eval\_point})\bigr)$ -- the
multilinear evaluation of the \emph{leaf} fractions at the accumulated
point.  The caller closes the loop by checking this against the
fingerprints $(\text{multiplicity}_i,\; \alpha - \text{fingerprint}_i)$
reconstructed from the main-trace PCS opening at \texttt{eval\_point}.

\subsection{Soundness Tests}

Six unit tests validate the protocol:

\begin{table}[h]
\centering
\begin{tabular}{ll}
\toprule
\textbf{Test} & \textbf{What it establishes} \\
\midrule
\texttt{fraction\_combine\_is\_associative\_on\_small\_input} & Algebraic sanity of $\oplus$ \\
\texttt{sum\_of\_balanced\_fractions\_has\_zero\_numerator}   & Target identity: balanced sends/receives $\Rightarrow N_m = 0$ \\
\texttt{logup\_gkr\_prove\_verify\_roundtrip\_on\_balanced\_input} & Honest prover round-trip \\
\texttt{logup\_gkr\_leaf\_claim\_matches\_mle\_evaluation}    & Reduction identity: final claim = $N_0(\mathsf{eval\_point})$ \\
\texttt{logup\_gkr\_rejects\_tampered\_combine\_identity}     & Soundness: corrupted combine round is rejected \\
\texttt{logup\_gkr\_rejects\_tampered\_root}                  & Soundness: corrupted root is rejected \\
\bottomrule
\end{tabular}
\caption{LogUp-GKR unit tests, all passing.}
\label{tab:logup-gkr-tests}
\end{table}

The \texttt{leaf\_claim\_matches\_mle\_evaluation} test is the key
soundness property: without it, the protocol might appear to verify but
would not actually bind the prover to the committed fraction leaves.

\subsection{Public API}

\begin{verbatim}
pub struct Fraction<EF>      { pub num: EF, pub denom: EF }
pub struct LogUpGkrProof<EF> {
    pub root: (EF, EF),
    pub rounds: Vec<(EF, EF, EF, EF)>,  // (n0, n1, d0, d1) per layer
    pub eval_point: Vec<EF>,
    pub leaf_claim: (EF, EF),
}

pub fn build_fraction_tree<EF>(leaves: &[Fraction<EF>]) -> Vec<Vec<Fraction<EF>>>;
pub fn sum_of_fractions<EF>(leaves: &[Fraction<EF>]) -> Fraction<EF>;
pub fn eval_mle<EF>(table: &[EF], r: &[EF]) -> EF;
pub fn eval_mle_first_var<EF>(table: &[EF], r: &[EF]) -> EF;

pub fn prove_logup_gkr<F, EF, Challenger>(
    leaves: &[Fraction<EF>], challenger: &mut Challenger
) -> LogUpGkrProof<EF>;

pub fn verify_logup_gkr<F, EF, Challenger>(
    proof: &LogUpGkrProof<EF>, challenger: &mut Challenger
) -> Option<(Vec<EF>, Fraction<EF>)>;
\end{verbatim}

The \texttt{eval\_point} returned is in MSB-first order (matching
\texttt{eval\_mle\_first\_var}), so the caller can evaluate the leaf MLE
directly at that point with no re-ordering.

\subsection{Integration Path (Phase 2b Steps 2 -- 4)}

Outstanding work to connect this protocol to the native prover:

\begin{enumerate}
  \item \textbf{Build the leaves from lookup events}. For each chip:
        \begin{itemize}
          \item Enumerate send/receive events;
          \item For each event, compute the fingerprint
                $f = \beta \cdot \text{payload}_0 + \beta^2 \cdot \text{payload}_1 + \dots$;
          \item Leaf numerator $= \pm \text{multiplicity}$
                ($+$ for sends, $-$ for receives);
          \item Leaf denominator $= \alpha - f$.
        \end{itemize}
  \item \textbf{Attach the GKR proof to \texttt{ShardProof}}.  Add a new
        optional field
        \texttt{logup\_gkr\_proofs: Option<Vec<LogUpGkrProof<Challenge<SC>>>>}
        with \texttt{\#[serde(default)]}.
  \item \textbf{Remove the permutation trace in WHIR mode}.  Skip
        \texttt{generate\_permutation\_trace} when LogUp-GKR is enabled.
        Drop the permutation-commit round from the Fiat-Shamir transcript.
  \item \textbf{Verifier: close the loop}.  After verifying the GKR layers,
        check the leaf claim against fingerprints recomputed from the
        opened trace values at \texttt{eval\_point}.
  \item \textbf{Unblock zerocheck}.  Once the permutation argument is
        gone, every chip becomes zerocheck-eligible; the Phase~2a
        limitation at \texttt{chip.permutation\_width() > 0} can be
        dropped.
\end{enumerate}

\subsection{Combined Soundness Picture}

\begin{table}[h]
\centering
\small
\begin{tabular}{lcccc}
\toprule
\textbf{Component} & \textbf{FRI} & \textbf{WHIR (pre-2a)} & \textbf{WHIR (2a only)} & \textbf{WHIR (2a + 2b + multi-point open)} \\
\midrule
Main trace binding    & \checkmark & \checkmark & \checkmark & \checkmark \\
Lookup integrity      & \checkmark & $\bigtriangleup$\footnotemark[1] & $\bigtriangleup$\footnotemark[1] & \checkmark (GKR) \\
Transition constraints & \checkmark & $\times$ & $\bigtriangleup$\footnotemark[2] & \checkmark (sumcheck) \\
Production-sound?     & \checkmark & $\times$ & $\times$ & \checkmark \\
\bottomrule
\end{tabular}
\caption{Soundness status across phases.}
\label{tab:soundness-phases}
\end{table}
\footnotetext[1]{Only the cumulative-sum balance is checked.}
\footnotetext[2]{Sumcheck transcript is bound, but the final claim is not yet checked against a PCS opening.}
